\documentclass[a4paper,12pt]{jsarticle}
\usepackage[margin=24mm]{geometry}
\usepackage{amsmath}
\usepackage{array}
\usepackage{booktabs}
\usepackage{enumitem}
\usepackage{longtable}
\usepackage{tabularx}

\setlength{\parindent}{0pt}
\setlength{\parskip}{0.7em}
\renewcommand{\arraystretch}{1.2}

\begin{document}

\begin{center}
{\LARGE 数学1A（学門ABDE） 授業紹介・シラバス概要}\\[0.4em]
2026年度春学期
\end{center}

\begin{tabularx}{\linewidth}{@{}>{\raggedright\arraybackslash}p{0.24\linewidth}X@{}}
\toprule
科目名 & 数学1A（学門ABDE）／（旧学門1345） \\
サブタイトル & 微分法の基礎と応用 \\
授業CD & 4317 \\
学期 & 2026年度春学期 \\
曜日時限 & 水曜1限 \\
キャンパス & 日吉 \\
教室 & 日吉23 \\
授業形態 & 対面授業（主として対面授業） \\
担当教員 & 中島秀太 \\
\bottomrule
\end{tabularx}

\section*{共通シラバスの概要}

\subsection*{授業科目の内容・目的・方法・到達目標}

１変数関数および多変数関数の微分法に関する基礎概念の理解と、それに裏打ちされた計算力を身につけることが本講義の目的です。講義の具体的内容は授業計画欄に記載されています。課題、レポート、小テスト、中間試験等の内容や実施時期および演習の方法は各担当者が判断します。なお、講義の内容とその順序は予告なしに変更される場合があります。

\subsection*{授業計画}

\begin{longtable}{@{}p{0.11\linewidth}p{0.83\linewidth}@{}}
\toprule
回 & 項目 \\
\midrule
\endhead
\toprule
回 & 項目 \\
\midrule
\endhead
\bottomrule
\endfoot
第1回 & １変数関数の極限と連続性 \\
第2回 & １変数関数の可微分性 \\
第3回 & 平均値の定理、ロピタルの定理 \\
第4回 & １変数関数のテイラー展開 \\
第5回 & テイラー展開の応用 \\
第6回 & 多変数関数の連続性、偏微分 \\
第7回 & 多変数関数の可微分性、合成関数の微分 \\
第8回 & 中間試験等 \\
第9回 & 陰関数定理 \\
第10回 & 陰関数の微分 \\
第11回 & 多変数関数のテイラー展開 \\
第12回 & 多変数関数の極値問題 \\
第13回 & ラグランジュの乗数法 \\
第14回 & まとめ \\
第15回 & 期末試験 \\
\end{longtable}


\subsection*{教科書}

慶應義塾大学理工学部数理科学科 編『数学1A・1B』を用います。生協で購入可能です。

\subsection*{担当教員から履修者へのコメント}

\begin{itemize}[leftmargin=2em]
\item 期末試験と同様に、課題等にも真摯に取り組んでください。
\item 授業は対面実施を基本とします。
\end{itemize}

\subsection*{質問・相談}

\begin{itemize}[leftmargin=2em]
\item 質問、相談の方法は各クラスの担当者が決めるので、各担当者に確認してください。
\item チュートリアルアワーという制度があり、大学院生の方々が質問を受け付けてくれます。
\end{itemize}

\section*{このクラスでの補足}

\subsection*{授業の進め方と準備学修}

\begin{itemize}[leftmargin=2em]
\item 小テストは原則として毎週 9:00--9:15 に実施します。
\item 小テストでは主として前回までの授業内容の確認を行います。
\item 授業では定義、定理、典型例を板書中心で説明し、必要に応じて演習問題を扱います。
\item 毎回の授業に向けて、教科書や配布資料の該当部分に事前に目を通してください。
\item 授業後は定義や定理の確認、例題の解き直し、小テスト範囲の復習を行ってください。
\end{itemize}

\subsection*{このクラスの成績評価}

このクラスでは、共通シラバスの範囲内で、成績を次の二つの計算結果のうち高い方で評価します。
\[
\max\left\{
\begin{array}{l}
\text{小テスト}(20\%) + \text{期末テスト}(80\%), \\[0.3em]
\text{小テスト}(10\%) + \text{中間テスト}(30\%) + \text{期末テスト}(60\%)
\end{array}
\right\}
\]

中間テストは 2026 年 6 月 3 日（水）に実施する予定です。詳細は K-LMS で案内します。

\subsection*{連絡と欠席}

\begin{itemize}[leftmargin=2em]
\item 連絡は K-LMS を通してお願いします。必要に応じて、対面でも相談を受けます。
\item 学期中の欠席が 3 回以下であれば、成績への影響はほとんどないので、報告は不要です。
\item 4 回以上欠席する見込みがある場合は、できるだけ早めに連絡してください。
\item 数回の欠席でも気になることがあれば、遠慮なく連絡してください。
\end{itemize}

\subsection*{春学期チュートリアルアワー（数学）}

\begin{itemize}[leftmargin=2em]
\item 時間は月曜日 13:30--15:30 です。
\item 2026 年 4 月 20 日（月）から開始します。
\item 場所は独立館地下1階 DB101, DB102 です。
\end{itemize}

\end{document}
